\section{The signed-excess estimate}\label{sec:signed-statement}

We now state the estimate that drives the proof. The definitions of
$M,B,D,L,\Lambda$, and $\tau$ are those of \Cref{sec:local}; in particular,
$L=2M-D$ and $\tau=|M|+D+\Lambda$ are measures on the base. In the curved case
the reference graph has small gradient, a fixed normal tube, and fixed bounds
through a sufficiently large fixed finite order $J$, as specified in
\Cref{sec:local}. All applications below use $C^J$ reference graphs
with these uniform bounds.

For a positive integer $j$, let $\exp_j$ and $\log_j$ denote $j$-fold
composition of the exponential and logarithm. For large $T$ set
\begin{equation}\label{eq:signed-rate}
 F_Q(T)=\exp_{2Q}\!\left(\sqrt{\log_{2Q}T}\right).
\end{equation}
Only sufficiently large arguments are used. Thus
\begin{equation}\label{eq:signed-rate-growth}
 (\log T)^p=o(F_Q(T))\quad(p<\infty),
 \qquad F_Q(T)=T^{o(1)}.
\end{equation}

\begin{theorem}[Signed excess]\label{sx:signed}
Fix a positive integer $Q$ and $\kappa>0$. Work in the interior normal-cylinder setting of
\Cref{sec:local}, with base coordinates scaled by a physical radius $l$ so
that $B_4$ is available. Suppose
\begin{equation}\label{eq:signed-smallness}
 \tau(B_4)\leq d=e^{-T},
 \qquad l+\sup|z_{\mathrm{phys}}|\leq e^{-\kappa T}.
\end{equation}
The supremum is over the restricted support above the available chart.
The statement also holds after an interior rescaling in the inherited
frames of \Cref{lt:heighttests}, provided in the curved case the second
smallness condition is imposed on the inherited parameter $\gamma$ from
\eqref{eq:lt-small-geometry}.
There are a constant $C_Q$ and a threshold $T_0$, depending only on the stated
fixed data, such that for $T\geq T_0$,
\begin{equation}\label{eq:signed-conclusion}
 \int w\,(L-C_Q\Lambda)\leq e^{-F_Q(T)}d
\end{equation}
whenever $w$ is supported in $B_1$, $0\leq w\leq1$, and $\Lip w\leq1$.
For a flat reference plane, take $\Lambda=0$ and omit the second condition
in \eqref{eq:signed-smallness}; constant sheet offsets need not be bounded.
\end{theorem}

The integer $Q$ is the reference count subtracted in $M$; the theorem
does not assume the pointwise bound $\theta\le Q$. The proof occupies the
next three sections and proceeds by induction on this reference count.
At the step $Q>1$, the hypothesis is the theorem for every integer
$1\le q<Q$.
If the asserted bound fails, \Cref{sec:contact} subtracts a small margin
from the weighted signed source and assigns its remaining positive mass
to selected centers and scales. Let $h$ denote the positive spatial
projection of this contact measure. At a contact, the smoothed difference
between the source and $h$ is nonnegative.
At larger scales only its cumulative integral with weight
$v^2\,\dd\log v$ is nonnegative. This remains true for contacts on
vertical sides of the stopping boundary. The construction itself uses
no smaller-integer assertion.

Induction enters through a configuration of $Q$ heights. If its entries
split into separated nonempty groups, sufficiently small working balls
contain actual stationary restrictions with positive reference counts
$q<Q$. The separated estimate in \Cref{sec:contact} applies
\Cref{sx:signed} at every such lower count and also gains a small positive
multiple of the tilt far from the configuration. In the main induction
step, \Cref{sec:excess} retains contacts whose outward vertical contains
a threshold equality and constructs an envelope of future signed mass.
Along the selected path beginning at that equality, the cumulative sign
makes the envelope control unsigned averages at larger radii up to the
chosen outer scale, with a fixed power of the radius ratio. This permits
a comparison with squared distance from finitely many height centers.
The separated gain pays for the regions where that distance function
changes between centers. The centers may merge, but at most $Q-1$ times,
so one of $Q$ fixed ranges of excess levels is traversed without a merger.

\Cref{sec:packing} extracts many scales with a definite normal flux on
that traversal. The higher height moment prevents this flux from lying
entirely where the excess measure has large density. A dyadic packing
argument then makes the union carrying all these scales too small to
support the original contact mass, closing the induction at $Q$.
For $Q=1$ every step involving separated groups or several wells is omitted.

\paragraph{Order of choices.}
Dimensions, a bound on $Q$, coordinate geometry, and $\kappa$ are fixed first.
Interior radii, covering constants, moment exponents, and the finite number
of derivative bounds are then fixed. Polynomial error exponents may be made
large enough for these data. For $Q>1$, the separated estimate uses a function
$b=b_Q$ with
\begin{equation}\label{eq:choice-b}
 b(T)\gg(\log T)^{200},\qquad
 F_{Q-1}(b(T)/2)\gg(\log T)^2,\qquad b(T)=T^{o(1)}.
\end{equation}
For example, a sufficiently large fixed multiple of
\[
 1+(\log T)^{201}
   +F_{Q-1}^{-1}((\log T)^3)
\]
suffices. Fixed increases of this choice will be made before fixing $T_0$.
The separation factors are exponentials of fixed multiples of $b$.
\Cref{ex:bands} verifies the additional rate comparisons required by the
finitely many mergers. Finally $T$ is taken large. A fixed weakening of
$\kappa$ changes only the threshold; this permits all the rescaled
lower-integer applications.

\paragraph{Margins and tails.}
Whenever a compactly supported chart is replaced by a Gaussian convolution,
the source includes a negative margin proportional to
$\eta(\tau+d\,\dd y)$, where $\eta>0$ is a chosen small parameter.
Contact scales lie below a chosen outer scale
$\rho$, with $|\log\rho|=o(T)$ in the main application. On the continuation
region, bounds at every larger scale control the Gaussian tails.
\Cref{ct:localization} gives the precise localization statement. Curved
first-variation errors contain the factor $e^{-\kappa T}$, which is smaller
than the margins used below. These conventions are retained at recursive
nodes, with lengths normalized by the node radius.
